\section{A deflated majorant in nine positive variables}
\label{sec:deflated-nine}

This contribution keeps the source and free-coordinate compiler of
\path{59-degree-separated/nine_free_adjunct} at ProveIt commit
\texttt{beb0192cb8dc36f78c369c080f2adfbc1887a486}.
It changes the polynomial that strictly majorizes a radical in the last
norm equation.  The resulting literal polynomial has nine strictly positive
unknowns, exact total degree \(1{,}573{,}266\), and a division-free evaluator
with \(504\) binary operations.  The predecessor has degree
\(1{,}910{,}404\) and \(473\) operations: the degree decreases by
\(337{,}138\), or approximately \(17.6475\%\), at the cost of \(31\)
additional operations.  Thus this is a degree/size tradeoff.

\subsection{New algebraic ingredient}

Use the inherited registers
\begin{align*}
 J&=4A-5,&D&=(A^2-1)C^2+1,&e&=iJC,\\
 E&=eDC,&F&=(A^2-1)E^2+1,&L_m&=CD+1-A,\\
 G&=A+FL_m,&H&=C+F(2R+1+(j-1)C),&I&=(G^2-1)H^2+1.
\end{align*}
The required domain is \(A\geq3\), \(C,i\geq1\).  Set
\begin{align*}
 N_*&=2G+L_m,&N_1&=2e^2D(D-1)L_m+3C,\\
 N_2&=2e^2(D-1)L_m,&N_3&=2e^2(L_m-C),&N_4&=2e^2C.
\end{align*}
The new polynomial is
\begin{align}
 T_7={}&e\bigl[(2048D^2-1024D-256)N_*
       -128N_1-80N_2-56N_3-42N_4 \notag\\
       &\hspace{22mm}+i^2(1568-448A+288C)\bigr]+640i+4096.
       \label{eq:deflated-T7}
\end{align}
No displayed register adds a quantified unknown.  They are abbreviations
for ordinary integer polynomials in the nine final inputs.

The simplifying relation is
\[
 e^2=16i^2D-i^2\bigl((40A-41)C^2+16\bigr).
\]
It cancels a complete high-degree part of the quotient remainder used by
the earlier majorant.  To make this exact, define
\begin{align*}
 U_6={}&e\bigl[(1024D^2-512D-128)N_*
             -64N_1-40N_2-28N_3-21N_4\bigr],\\
 T_6={}&U_6-224i^2e(A-1-3C)+560i^2e+320i+2048.
\end{align*}
Then \(T_7=2T_6-1056i^2eC\).  The variable-dependent subtraction is
justified by the following exact certificate.

Put
\begin{align*}
 t_7={}&2048D^7-1024D^6-256D^5-128D^4
           -80D^3-56D^2-42D-33,\\
 d_7={}&109824D^6+36608D^5+18304D^4+9984D^3
           +5460D^2+2772D+1089,\\
 R_7={}&D^5T_7-et_7N_*.
\end{align*}
Exact expansion gives
\[
 t_7^2-2048^2D^{13}(D-1)=d_7>0.
\]
Since the sum of the absolute negative coefficients of \(t_7\) is
\(1619<2048\), one has \(t_7>0\) for \(D\geq1\).  The delivered
file \texttt{data/order7\_majorant\_positive\_certificate.json} gives
the complete exact identity
\begin{equation}
 R_7-4096D^5
   =\sum_{(u,v,w)\in\mathcal S} a_{uvw}(A-3)^u(C-1)^v i^w,
 \qquad |\mathcal S|=264,
 \label{eq:deflated-positive-certificate}
\end{equation}
where every \(a_{uvw}\) is a strictly positive integer.  At \(A=3,C=1\)
the right side is
\(49{,}748{,}823i+293{,}313{,}636i^3>0\).
The checker reconstructs every coefficient from
\eqref{eq:deflated-T7}; this is a finite polynomial identity and a uniform
positivity proof, rather than a test on sampled parameter values.

For completeness, write \(X=e^2D^3(D-1)\).  Then
\(DF=X+D\) and \(G=(X/D)L_m+CD+1\).  The elementary inequality
\(\sqrt{X+D}<\sqrt X+D/(2\sqrt X)\) gives
\[
 4096\sqrt{DF}\,G
 <2048eD^2\sqrt{1-1/D}\,N_*+4096v,
 \qquad
 v=\frac{CD+1}{2eD\sqrt{1-1/D}}<1.
\]
The final inequality follows from \(D\geq A^2\geq9\) and \(e\geq7C\).
Combining it with \eqref{eq:deflated-positive-certificate} and the square
identity proves the unconditional bound
\begin{equation}
             T_7>4096\sqrt{DF}\,G>0.
             \label{eq:deflated-majorant}
\end{equation}
The divisions occur only in this proof, not in the polynomial.

\subsection{Final polynomial and semantic equivalence}

Keep the inherited positive source parameters and normalized margin:
\begin{align*}
 B&=2^{2134939}b^2,&N&=2^{17079501}bB^{1201},\\
 t&=bB^{116},&m&=K^2(4t-g)-16t(C-KY)^2,\\
 \Omega&=Jm,&E_0&=J\bigl(n+(bw)^2\bigr)+Q.
\end{align*}
Set
\[
\begin{aligned}
 a_0&=(P^2-1)K^2+1,& a_1&=4096^2DFI,\\
 V_0&=PK,& V_1&=T_7H,\\
 \Theta&=E_0-\Omega(V_0+V_1).
\end{aligned}
\]
The equation is
\begin{equation}
 F_9=\bigl(\Theta^2-\Omega^2(a_0+a_1)\bigr)^2
            -4\Omega^4a_0a_1=0.
 \label{eq:deflated-final}
\end{equation}
Its only unknowns are \(b,g,h,s,w,C,i,j,n\), all strictly positive.
The source coefficients, source positions and ratio conditions are exactly
those of the pinned predecessor.  The four full native coefficient
payloads are supplied; there is no unspecified program-index parameter.

Equation \eqref{eq:deflated-majorant} gives \(V_1>\sqrt{a_1}\), and
\(V_0^2-a_0=K^2-1>0\).  Also
\(E_0=Jn+(J-2)(bw)^2+3bwC+2>0\).  At a zero, \(m=0\) is impossible.
Real factorization yields roots \(r_0^2=a_0,r_1^2=a_1\) with
\[
 n+(bw)^2+Q/J=m(V_0+V_1+r_0+r_1).
\]
Both the left side and the right-hand parenthesis are strictly positive.
Consequently \(m>0\) follows before any integrality assertion.  The
inherited margin argument recovers \(C>2R+1+bw\) and \(C>K\), hence
\(a_1>a_0\).  The rational sum \(r_0+r_1\) is therefore nonzero, so
\(2(r_0+r_1)r_0=(r_0+r_1)^2+a_0-a_1\) makes the roots rational and
then integral.  The same rational-square argument makes \(r_1/4096\)
integral, and the factor equation gives \(J\mid Q\).  These are the
unchanged ratio hypotheses.

Conversely, take the inherited ratio witnesses and positive roots, and put
\(n=m(V_0+V_1+r_0+r_1)-(bw)^2-Q/J\).  It is an integer, and it is
positive because \((bw)^2+Q/J<2D\), \(m\geq1\), and
\(V_1>4096D\).  It makes \eqref{eq:deflated-final} vanish.  Thus the
positive zero set is equivalent to the pinned source zero set.  Composing
with its native compiler and K215 calibration gives, in ZFC, absence of
positive zeros if and only if \(\operatorname{Con}(\mathrm{ZFC})\).
This is an ordinary mathematical proof with named inherited dependencies;
the new semantic equivalence is not claimed to have been kernel checked.

\subsection{Exact degree and finite verification}

Let \(\delta_j=V_j^2-a_j\) and
\(B_*=E_0^2-2E_0\Omega(V_0+V_1)+\Omega^2(\delta_0+\delta_1)\).
The two-defect identity gives
\[
 F_9=B_*^2+4\Omega^2V_0V_1B_*
          +4\Omega^4(\delta_0V_1^2+a_0\delta_1).
\]
The exact majorant defect identity is
\begin{align*}
 D^{10}(T_7^2-4096^2DFG^2)
  ={}&e^2d_7N_*^2
     +2048^2D^{10}\bigl[XL_m^2-4DG(CD+1)\bigr]\\
    &+2et_7N_*R_7+R_7^2.
\end{align*}
Here \(\deg R_7=187305\), with highest weighted term
\(2880A^{10}C^{12}i^3\).  The mixed term uniquely dominates the last
identity.  It yields
\[
 \deg(T_7^2-4096^2DFG^2)=280959,\qquad
 \deg\delta_1=598873,\qquad \deg B_*=786633.
\]
The square \(B_*^2\) uniquely dominates the final two-defect expansion,
proving \(\deg F_9=1{,}573{,}266\).  A separate circuit for
\(D^{20}F_9\) supplies an exact nonzero diagonal leading coefficient
and the matching total-degree upper bound.  Its extra operations are not
included in the primary evaluator count.

Running \texttt{python3 verify.py} checks all four native byte payloads,
the source isolation data, the 264-term positivity identity, the defect
identity, every delivered circuit against its deterministic builder,
graph liveness, exact degrees, costs and domain translation.  It also
performs 480 comparisons with a separately written literal evaluator.
The check uses only Python's standard library.  These finite checks support
the implementation and identities; the unbounded semantic endpoint uses
the explicitly credited source and compiler theorems.
